\paragraph{A precise perfect-test relative-cohomology theorem.}
The following gives a modern derived construction directly relevant to
the relative-cohomology programme on cote~43, folios~95--96.
Its sites and comparison hypotheses are specified independently; we
do not identify them with an incompletely reconstructed historical site.

\begin{theorem}[Known: perfect Greenberg realization and relative cohomology]\label{B:perfect-Greenberg-cohomology}
Let $R$ be a complete discrete valuation ring with fraction field $K$,
uniformizer $\pi$, and perfect residue field $k$ of characteristic $p>0$.
Equal and mixed characteristic are both allowed. Fix a coefficient map
$W(k)\to R$. For a finite product $L$ of perfect field extensions of $k$,
put
\[
 R(L)=\varprojlim_n
 \bigl(W_n(L)\otimes_{W_n(k)}R/\mathfrak m_R^n\bigr),
 \qquad K(L)=R(L)[1/\pi].
\]
For $R=k[[t]]$ with its specified coefficient field, these are $L[[t]]$
and $L((t))$. For $L$ a field, $R(L)$ is a complete discrete valuation
ring with residue field $L$.

Let $S_{\mathrm{pa}}$ be the étale site on finite products of perfect
field extensions of $k$, $S_{\mathrm{pf}}$ the pro-fppf site on
quasi-compact quasi-separated perfect $k$-schemes, and
$S_{\mathrm{ir}}$ the pro-étale site on ind-rational $k$-algebras.
Here rational algebras are finite products of perfections of finitely
generated field extensions, and ind-rational algebras are filtered
unions of rational subalgebras. In the pro-fppf site, coverings are
finite jointly surjective families flat of profinite presentation in
the perfect sense. The inclusions induce premorphisms
\[
 h:S_{\mathrm{pf}}\to S_{\mathrm{pa}},\qquad
 \widetilde h:S_{\mathrm{pf}}\to S_{\mathrm{ir}},\qquad
 \mathsf a=R\widetilde h_*Lh^*:
 D(S_{\mathrm{pa}},\mathbf{Ab})\to D(S_{\mathrm{ir}},\mathbf{Ab}).
\]
These premorphisms are not assumed to have exact inverse image on
arbitrary sheaves of sets.
\begin{enumerate}
\item Let $\mathcal P'_{\mathrm{fc}}$ consist of commutative
proalgebraic groups representable as extensions of perfections of
abelian varieties by perfect affine groups, with finite étale component
group. Restriction to $S_{\mathrm{pa}}$ is fully faithful on this
category, and its derived counit is an isomorphism:
\[
 Lh^*(P|_{S_{\mathrm{pa}}})\simeq P
 \quad\text{in }D(S_{\mathrm{pf}},\mathbf{Ab}).
\]
\item The functors $R(-),K(-)$ induce premorphisms
$\pi_R:R_{\mathrm{fppf}}\to S_{\mathrm{pa}}$ and
$\pi_K:K_{\mathrm{fppf}}\to S_{\mathrm{pa}}$. Define
\[
 \mathcal R_R(F)=\mathsf aR\pi_{R,*}F,\qquad
 \mathcal R_K(F)=\mathsf aR\pi_{K,*}F,\qquad
 \mathcal R_x(F)=
 \operatorname{fib}\bigl(\mathcal R_R(F)\to\mathcal R_K(F|_K)\bigr).
\]
They give a functorial localization triangle.
Call $Q\in D(S_{\mathrm{pa}},\mathbf{Ab})$ $h$-acyclic if its unit
$Q\to Rh_*Lh^*Q$ is an isomorphism. If $R\pi_{R,*}F$ is
$h$-acyclic, then
\[
 R\Gamma(L_{\mathrm{pro\acute et}},\mathcal R_R(F))
 \simeq R\Gamma(R(L)_{\mathrm{fppf}},F).
\]
The analogous assertion holds over $K$. If both unit conditions hold
for $F$ and $F|_K$, the analogous comparison holds with support on
the closed fibre. Under this unit hypothesis, if both $F$ and
$\mathcal R_R(F)$ are bounded below, there is the spectral sequence
\[
 E_2^{ij}=
 H^i(L_{\mathrm{pro\acute et}},\mathcal H^j\mathcal R_R(F))
 \Longrightarrow H^{i+j}(R(L)_{\mathrm{fppf}},F).
\]
For algebraically closed $L$, this reduces to
$\mathcal H^q\mathcal R_R(F)(L)=H^q(R(L)_{\mathrm{fppf}},F)$.
\item If $G/R$ is a smooth commutative group scheme of finite type,
write $G_n=G\times_RR/\mathfrak m_R^n$ and
\[
 \mathbf G=\varprojlim_n
 \bigl(\operatorname{Gr}_{R/\mathfrak m_R^n}(G_n)\bigr)^{\mathrm{perf}}.
\]
Then $\mathbf G\in\mathcal P'_{\mathrm{fc}}$ and
\[
 0\longrightarrow U_G\longrightarrow\mathbf G
 \longrightarrow G_k^{\mathrm{perf}}\longrightarrow0
\]
is exact, with $U_G$ connected pro-unipotent.
Consequently $\pi_0(\mathbf G)=\pi_0(G_k)$, and
$\mathcal R_R(G)\simeq\mathbf G[0]$.
The unit condition in the preceding clause holds.
\item If $N/R$ is finite flat and commutative, then
$\mathcal H^0\mathcal R_R(N)$ is finite étale and agrees with
$\mathcal H^0\mathcal R_K(N_K)$,
$\mathcal H^1\mathcal R_R(N)$ is connected pro-unipotent, and
$\mathcal H^q\mathcal R_R(N)=0$ for $q\geq2$.
The object $R\pi_{R,*}N$ is $h$-acyclic.
\end{enumerate}
\end{theorem}

\begin{proof}[Derivation from the opened sources]
The sites, ring functors and premorphisms are those of
\cite[Definitions~3.2, 3.4, 3.7--3.8, 7.4--7.7, 8.1--8.4;
Proposition~8.3]{BSuzuki2021}; all locators here refer to its opened
7~October~2019 author version. Finite étale residue-field extensions
lift to finite étale extensions of Witt rings. Tensoring and completing
with $R$ preserves these finite free extensions and their covering
fibre products. This proves the asserted premorphism property, rather
than asserting exactness of an arbitrary inverse-image functor.

For the first clause, the specified category is Definitions~3.12--3.13
of that source. Theorem~3.15 applies to the larger category $E_k$, and
therefore gives the displayed derived counit. Its degree-zero part
and adjunction give, for $P,Q\in\mathcal P'_{\mathrm{fc}}$,
\[
 \operatorname{Hom}_{S_{\mathrm{pa}}}(P,Q)
 =\operatorname{Hom}_{S_{\mathrm{pf}}}(h^*P,Q)
 =\operatorname{Hom}_{S_{\mathrm{pf}}}(P,Q).
\]
This is precisely Proposition~7.1, or Theorem~1.1.
The source proves the derived assertion in its complete
Sections~4--6 using the Mac Lane resolution and pro-fppf covers of
generic points. Proposition~3.9 explicitly rules out replacing that
argument by exactness on arbitrary sheaves of sets.

For the second clause, $\mathsf a$ and the relative functors are
Definitions~7.7 and~8.4. Derived functors preserve the defining fibre
triangle, giving (8.1). Factor
$h=\alpha\epsilon\widetilde h$ as in Definition~7.5, with
$\epsilon$ the change from pro-étale to étale topology on ind-rational
algebras and $\alpha_*$ exact. The \emph{unit} hypothesis gives
\[
 R\pi_{R,*}F
 \simeq Rh_*Lh^*R\pi_{R,*}F
 \simeq\alpha_*R\epsilon_*\mathcal R_R(F).
\]
Derived sections at $L$ give the claimed comparison, exactly as in
Proposition~8.5. The displayed spectral sequence is the bounded-below hypercohomology
spectral sequence of the relative object. We explicitly retain a lower
bound on that object; no general t-exactness of $Lh^*$ is inferred.
This extra bound is automatic in the third and fourth clauses. Pro-étale covers of an algebraically
closed field split, so sections there are exact. Applying the argument
over $K$ and then taking the localization fibre gives the supported
version. The counit of the first clause and this unit are different
maps; the second assertion does not discard its acyclicity hypothesis.

For the third clause, finite Greenberg representability identifies
the points of its $n$th term on perfect residue-field extensions with
$G(R(L)/\mathfrak m_R^n)$
\cite[Proposition--Definition~7.10, (7.16)]{BBertapelleGonzalezAviles2022}.
Smoothness gives surjectivity of reduction. The successive transition
kernels become vector groups after perfection, so the kernel of the
map to $G_k^{\mathrm{perf}}$ is connected pro-unipotent. The first
stage has finite component group because $G_k$ is of finite type.
These are exactly the argument and conclusions of
\cite[Proposition~(3.4.2)(a), pp.~48--49]{BSuzuki2020},
whose locators refer to the opened 22~March~2018 author version.
The infinite point formula follows by taking the complete compatible
system of points; the general algebraization justification is
\cite[Proposition~18.2]{BBertapelleGonzalezAviles2022}, using
\cite[Theorem~4.1, Remarks~4.3 and~4.6]{ROOT_BhattAlgebraization2014}.

Smooth fppf cohomology agrees with étale cohomology. The exactness
statement in \cite[Proposition~(3.4.1)]{BSuzuki2020} says that the
relative étale direct image has no higher cohomology: étale covers of
the complete valuation-ring tests are locally refined by lifted
residue-field covers. Thus the old relative object is
$\mathbf G[0]$. Propositions~10.8--10.9 of
\cite{BSuzuki2021} identify it with the new functor above and give
the unit condition. Proposition~7.10(b) also verifies that condition
for the tower with connected unipotent transition kernels. The
vanishing is a statement about relative cohomology \emph{sheaves};
ordinary cohomology for non-algebraically-closed $L$ is still computed
by the spectral sequence in the second clause.

For the fourth clause, the standard smooth resolution
\cite[Proposition~2.2.1, pp.~25--27]{BBegueri1980}, applied to the
Cartier dual of $N$, gives
\[
 0\longrightarrow N\longrightarrow H\longrightarrow J\longrightarrow0
\]
with $H,J$ smooth affine commutative and with connected fibres.
The first smooth group is a Weil restriction of $\mathbb G_m$ from
the finite flat Cartier dual; its geometric fibres are unit groups of
finite algebras and are connected. Its quotient by the finite
subgroup is connected as well. Applying the third clause gives
\[
 0\to\mathcal H^0\mathcal R_R(N)\to\mathbf H\to\mathbf J
 \to\mathcal H^1\mathcal R_R(N)\to0
\]
and vanishing in degrees at least two.
Proposition~(3.2.10) of \cite{BSuzuki2020} identifies the degree-zero
object with the finite étale group defined by $N(K^{\mathrm{ur}})$
and its residue-field Galois action. Since $N$ itself is the locally
finite part, the $R$ and $K$ point functors agree.
The cokernel is connected proalgebraic and is annihilated by the
order of $N$. Its semiabelian part is therefore zero, so it is
pro-unipotent. These are the complete steps of
\cite[Proposition~(3.4.2)(b)]{BSuzuki2020}. The comparison and the
unit condition are again Propositions~10.8--10.9 of \cite{BSuzuki2021}.
\end{proof}

\paragraph{The perfection qualification and the original programme.}
Bégueri's own Section~3.1 fixes mixed characteristic and an
algebraically closed residue field. Its Section~4.1 and the remarks
after Proposition~6.2.3, pp.~76--77, construct perfect Greenberg
realizations and a derived functor on groups of multiplicative type.
They provide a direct precursor for the present relative construction.
The nonperfect kernel formula of Lemma~4.1.1(2) is nevertheless false:
for $R=W(k)$ and $G=\mathbb G_a$, the proposed kernel map has kernel
$\alpha_p$
\cite[Remark~14.22]{BBertapelleGonzalezAviles2022}.
Perfection removes that defect, as Remark~15.9 explains and as
\cite[footnote~13, p.~48]{BSuzuki2020} explicitly requires.
We use perfect groups throughout. Neither this theorem nor the
finite-partition comparison identifies every historical test object
or proves the crossed-out injective-preservation argument on
cote~43/96. Those are scope limits on the historical interpretation,
not silently assumed steps in the proved modern theorem.
